\section{Quadratic relations and projective ramification}
\label{sec:relations-ramification}
All schemes in this section are over $\C$. Put
$N=e(e+1)/2$ and $p=N-e$. For $R\in\Gr(e,\Sym^2V)$ let
$\gamma_R:\Sym^2V\twoheadrightarrow S_R$ be the quotient map.
We regard elements of $\Sym^2V$ as homogeneous quadratic functions
on $V^*$. The algebra $B_R$ introduced above is local, with maximal
ideal $V\oplus S_R$, square $S_R$, and cube zero; hence its Hilbert
function is $(1,e,p)$.

Let $\mathcal H$ be the tautological hyperplane bundle on
$P=\Pj(V^*)$. Contraction is normalized by
$c_\alpha(vw)=\alpha(v)w+\alpha(w)v$.

\begin{lemma}[Restriction--contraction identity]
\label{lem:restriction-contraction}
There is an exact sequence of vector bundles on $P$
\begin{equation}\label{eq:contraction-sequence}
 0\longrightarrow\Sym^2\mathcal H
 \longrightarrow\Sym^2V\otimes\mathcal O_P
 \xrightarrow{\ c\ }V\otimes\mathcal O_P(1)
 \longrightarrow0.
\end{equation}
Consequently there is a natural isomorphism
\begin{equation}\label{eq:relation-cokernel}
 \operatorname{coker}\bigl(\Sym^2\mathcal H\to
                      S_R\otimes\mathcal O_P\bigr)
 \simeq
 \operatorname{coker}\bigl(R\otimes\mathcal O_P
                    \xrightarrow{c_R}V\otimes\mathcal O_P(1)\bigr).
\end{equation}
Both maps are square, of sizes $p$ and $e$. Their determinant ideals
are equal, and define a degree-$e$ hypersurface whenever the determinant
is not identically zero.
\end{lemma}
\begin{proof}
Locally choose $\alpha$ and $\ell\in V$ with $\alpha(\ell)=1$.
Then $V=H\oplus\C\ell$ and
\[
 \Sym^2V=\Sym^2H\oplus H\ell\oplus\C\ell^2.
\]
Contraction kills the first summand, maps $h\ell$ to $h$, and maps
$\ell^2$ to $2\ell$. This proves exactness, including local splitting;
the twist records rescaling of $\alpha$. Both cokernels in
\eqref{eq:relation-cokernel} are the quotient of
$\Sym^2V\otimes\mathcal O_P$ by
$R\otimes\mathcal O_P+\Sym^2\mathcal H$. Equality of their zeroth
Fitting ideals proves the scheme-theoretic determinant assertion.
Finally $\det c_R$ is a section of
$\det V\otimes(\det R)^*\otimes\mathcal O_P(e)$.
This also agrees with
$\det(\Sym^2\mathcal H)^*$, whose nonconstant factor is
$\mathcal O_P(e)$ because $\det\mathcal H=\mathcal O_P(-1)$ and
$\rank\mathcal H=e-1$.
\end{proof}

\begin{proposition}[The ramification interpretation]
\label{prop:quadratic-ramification}
Suppose $R$ is basepoint-free on $P$. Its evaluation morphism
$\phi_R:P\to\Pj(R^*)$ is finite, surjective, and has degree
$2^{e-1}$. The determinant scheme of
$\gamma_R|_{\Sym^2\mathcal H}$ is precisely its ramification scheme
\begin{equation}\label{eq:jacobian-ramification}
 Y_R=V\!\left(\det\left(\frac{\partial q_i}{\partial a_j}\right)_{i,j=1}^e\right)
       \subset\Pj(V^*),
\end{equation}
where $q_1,\ldots,q_e$ is any basis of $R$ and $a_1,\ldots,a_e$
are linear coordinates on $V^*$. Moreover
\begin{equation}\label{eq:basepoint-firstorder}
 R\text{ is basepoint-free}
 \quad\Longleftrightarrow\quad
 \gamma_R(HV)=S_R\text{ for every hyperplane }H\subset V.
\end{equation}
\end{proposition}
\begin{proof}
Basepoint-freeness gives a morphism with
$\phi_R^*\mathcal O(1)=\mathcal O_P(2)$. A positive-dimensional
projective fibre contains a curve; this would make $\mathcal O_P(2)$
trivial on that curve, contrary to ampleness. Thus the proper morphism
is quasi-finite and hence finite. Its image has dimension $e-1$ and
is the entire target. The intersection of $e-1$ pulled-back hyperplanes
has degree $2^{e-1}$, which is the degree of the finite morphism.
In characteristic zero it is generically separable.

The differential on the affine cones is the displayed Jacobian $J$,
which is the transpose of the contraction matrix in the chosen bases.
Euler's identity is $J(a)a=2(q_1(a),\ldots,q_e(a))$. On a source chart
$a_j\ne0$ and target chart $q_i\ne0$, choose the radial vectors as the
first vectors in the respective tangent frames. The induced map on
the radial line is multiplication by a unit. Passing to the quotient
by that line gives the differential of the projective morphism.
Expansion along this first column shows that its determinant and
$\det J$ differ by a unit. Equivalently their cokernels agree after
removing this invertible one-dimensional summand. This identifies the
ramification Fitting ideal, including its scheme structure.
Lemma~\ref{lem:restriction-contraction} identifies the other determinant.

For the last assertion, a functional on $S_R$ annihilating
$\gamma_R(HV)$ pulls back to a functional on $\Sym^2V$ annihilating
both $R$ and $HV$. If $H=\ker\alpha$, then
$(HV)^\perp=\C\alpha^2$. Thus failure of surjectivity is equivalent
to $q(\alpha)=0$ for every $q\in R$, exactly a base point of the
quadratic system. This proves both implications.
\end{proof}

\begin{proposition}[A nonempty open class with moduli]
\label{prop:admissible-ramification}
For $e=3$ or $e=4$, the set
\[
 G_e^\circ=\{R\in\Gr(e,\Sym^2V):
       R\text{ is basepoint-free and }Y_R\text{ is smooth}\}
\]
is nonempty, open, and invariant under $\operatorname{PGL}(V)$.
For every $R\in G_e^\circ$, the restriction
$\gamma_R|_{\Sym^2H}$ has rank $p-1$ on $Y_R$ and rank $p$
off $Y_R$. The variety $Y_R$ is smooth and integral. It has genus one
for $e=3$, and is a K3 surface for $e=4$.

Two algebras $B_R,B_{R'}$ are isomorphic if and only if $R,R'$ lie
in the same $\operatorname{PGL}(V)$ orbit. In particular
\begin{equation}\label{eq:relation-moduli-deficit}
 \dim G_e^\circ-\max_{R\in G_e^\circ}
          \dim\bigl(\operatorname{PGL}(V)\cdot R\bigr)
 \ \ge\ ep-e^2+1
 =\begin{cases}1,&e=3,\\9,&e=4.\end{cases}
\end{equation}
Here the left side is an orbit-dimension deficit; no fine moduli scheme
is being asserted.
\end{proposition}
\begin{proof}
We give a parameter-space proof of nonemptiness rather than infer
smoothness by specializing a primary decomposition. Fix a $p$-dimensional
space $S$ and let $T=\operatorname{Hom}(\Sym^2V,S)$. Over $T\times P$,
evaluation on $\Sym^2\mathcal H$ is a surjective vector-bundle map to
the bundle of $p$ by $p$ matrices. For a fixed $H$, corank at least two
therefore has codimension four in $T$. Its universal locus in
$T\times P$ has dimension $\dim T+e-1-4<\dim T$. Projection is
proper because $P$ is projective, so its image is a proper closed
subset of $T$.

Similarly $HV$ has dimension $p+e-1$. Restriction to it is a
surjective evaluation map to the space of $p$ by $(p+e-1)$ matrices.
Nonsurjectivity has codimension $e$, so the universal bad locus has
dimension $\dim T-1$. Its proper image is again a proper closed
subset. Remove both images, remove nonsurjective $\gamma$, and
require the restriction at one fixed hyperplane to be invertible.
These are nonempty open conditions in the irreducible affine space
$T$, so their intersection $T_1$ is nonempty. The last condition
ensures that the determinant is not the zero section.

On $T_1\times P$ the universal determinant hypersurface is smooth:
evaluation on the square matrix is a submersion, and the determinant
hypersurface is smooth at every rank-$p-1$ matrix. Every fibre is a
nonempty degree-$e$ hypersurface. The generic fibre over $T_1$ is regular,
since its local rings are localizations of regular rings, and is
geometrically regular in characteristic zero. It is therefore smooth.
The nonsmooth locus of the proper morphism to $T_1$ is closed; its
proper image misses the generic point. Deleting this image gives a
nonempty open $T_2$ with all determinant fibres smooth.

The map from the space of surjections $\Sym^2V\to S$ to
$\Gr(e,\Sym^2V)$ takes a tensor to its kernel and is the frame
bundle of the universal quotient. The preceding conditions depend
only on that kernel. Basepoint-freeness and smoothness of the
projective hypersurface are open conditions by the same proper-image
argument. Thus $G_e^\circ$ is open and contains the image of $T_2$;
it is invariant by its intrinsic definition. Smoothness of a
nonzero determinant hypersurface excludes matrices of corank at
least two, since their cofactors vanish and hence all derivatives
of the determinant in local base coordinates vanish.

A smooth hypersurface of positive degree in $\Pj^{e-1}$ is integral
here: for $e-1\ge2$, distinct positive-degree components intersect,
and their intersection would be singular. Adjunction gives
$\omega_{Y_R}\simeq\mathcal O_{Y_R}$. The hypersurface sequence
$0\to\mathcal O_P(-e)\to\mathcal O_P\to\mathcal O_{Y_R}\to0$
gives genus one when $e=3$, and $H^1(Y_R,\mathcal O_{Y_R})=0$ when
$e=4$. The latter is a K3 surface by the algebraic definition;
compare the standard quartic example in
\cite[Chapter 1, Example 1.3(i)]{HuyK3}.

An algebra isomorphism preserves the maximal ideal and its square.
On the associated graded algebra it induces an isomorphism of $V$
which takes the kernel of quadratic multiplication to the other
kernel. Conversely such a linear isomorphism induces a graded
algebra isomorphism. Scalars act trivially on relation subspaces,
so the equivalence is exactly the stated projective-linear orbit
relation. Finally $\dim\Gr(e,N)=e(N-e)=ep$ and every orbit has
dimension at most $e^2-1$. This proves the estimate and, in
particular, excludes a finite-orbit interpretation of the class.
\end{proof}

\begin{remark}\label{rem:ramification-scope}
The restrictions $e=3,4$ enter the nonemptiness proof at the
codimension-four corank-two incidence. For $e\ge5$ this dimension
argument no longer excludes its projection. The universal flag theorem
below handles additional restriction-rank strata, but the present
smooth-ramification theorem does not silently extend to them.
The construction of the finite quadratic morphism, the Jacobian
criterion and the quartic K3 calculation are classical ingredients.
They are used here to identify what the embedded part of the
multiplication-failure scheme intrinsically remembers.
\end{remark}
